% Generated by Cir2Tikz/v1rels.hs -- do not edit by hand.
% Build with:  pdflatex v1rels.tex   (run from the Cir2Tikz directory)
\documentclass[11pt]{article}
\usepackage[margin=1.4cm,landscape,a4paper]{geometry}
\usepackage{amsmath,amssymb,longtable,graphicx}
\usepackage{tikzit}
\input{circuits.tikzstyles}
\definecolor{normGreen}{RGB}{214,243,221}
\definecolor{normRed}{RGB}{249,196,196}
\definecolor{normYellow}{RGB}{253,246,208}
\pagestyle{empty}
\setlength{\LTpre}{0pt}\setlength{\LTpost}{0pt}
\setcounter{LTchunksize}{200}
\newlength{\figwidth}
\newsavebox{\figbox}
\newcommand{\fitfig}[1]{%
  \sbox{\figbox}{\tikzfig{#1}}%
  \ifdim\wd\figbox>\figwidth
    \sbox{\figbox}{\resizebox{\figwidth}{!}{\usebox{\figbox}}}%
  \fi
  $\vcenter{\hbox{\usebox{\figbox}}}$}
\begin{document}
\setlength{\figwidth}{\dimexpr\textwidth-9em\relax}
\begin{center}
  {\Large The relation set of \texttt{Simplified-V1.Syntactics}}\\[2pt]
  {\footnotesize \texttt{Clifford-Relations.Base.\_SRel,\_===\_} together
   with the structural rules of \texttt{Circuit.Base.Lift-Relation}.\\
   Diagrams read left to right in matrix-multiplication order (the
   \emph{rightmost} gate is applied first), so each transcribes its Agda
   word letter for letter.  Wires are numbered from the bottom: $w$ sits on
   wire $0$, $w\uparrow$ on wire $1$, and $w\downarrow$ is $w$ itself.}
\end{center}
\begin{longtable}{@{}r@{\quad}l@{}}
\multicolumn{2}{@{}l}{\bfseries Derived words (definitions)}\\*[2pt]
{\footnotesize $Z$} & \fitfig{v1/def-Z}\\[7pt]
{\footnotesize $X$} & \fitfig{v1/def-X}\\[7pt]
{\footnotesize $R$} & \fitfig{v1/def-R}\\[7pt]
{\footnotesize $M_x$} & \fitfig{v1/def-M}\\[7pt]
{\footnotesize $\bot\top$} & \fitfig{v1/def-bt}\\[7pt]
{\footnotesize $\top\bot$} & \fitfig{v1/def-tb}\\[7pt]
\multicolumn{2}{@{}l}{\bfseries One-wire axioms}\\*[2pt]
{\footnotesize order-S} & \fitfig{v1/order-S}\\[7pt]
{\footnotesize order-H} & \fitfig{v1/order-H}\\[7pt]
{\footnotesize M-power} & \fitfig{v1/M-power}\\[7pt]
{\footnotesize semi-MR} & \fitfig{v1/semi-MR}\\[7pt]
{\footnotesize order-SH} & \fitfig{v1/order-SH}\\[7pt]
{\footnotesize comm-HHSHHS} & \fitfig{v1/comm-HHSHHS}\\[7pt]
\multicolumn{2}{@{}l}{\bfseries Two-wire axioms}\\*[2pt]
{\footnotesize semi-M\ensuremath{\uparrow}CZ} & \fitfig{v1/semi-M-up-CZ}\\[7pt]
{\footnotesize semi-M\ensuremath{\downarrow}CZ} & \fitfig{v1/semi-M-dn-CZ}\\[7pt]
{\footnotesize rel-X\ensuremath{\uparrow}-CZ} & \fitfig{v1/rel-Xup-CZ}\\[7pt]
{\footnotesize rel-X\ensuremath{\downarrow}-CZ} & \fitfig{v1/rel-Xdn-CZ}\\[7pt]
{\footnotesize order-CZ} & \fitfig{v1/order-CZ}\\[7pt]
{\footnotesize comm-CZ-S\ensuremath{\downarrow}} & \fitfig{v1/comm-CZ-Sdn}\\[7pt]
{\footnotesize comm-CZ-S\ensuremath{\uparrow}} & \fitfig{v1/comm-CZ-Sup}\\[7pt]
{\footnotesize selinger-c10} & \fitfig{v1/selinger-c10}\\[7pt]
{\footnotesize selinger-c11} & \fitfig{v1/selinger-c11}\\[7pt]
\multicolumn{2}{@{}l}{\bfseries Three-wire axioms}\\*[2pt]
{\footnotesize selinger-c12} & \fitfig{v1/selinger-c12}\\[7pt]
{\footnotesize selinger-c13} & \fitfig{v1/selinger-c13}\\[7pt]
{\footnotesize selinger-c14} & \fitfig{v1/selinger-c14}\\[7pt]
{\footnotesize selinger-c15} & \fitfig{v1/selinger-c15}\\[7pt]
\multicolumn{2}{@{}l}{\bfseries Structural rules (Circuit.Base.Lift-Relation)}\\*[2pt]
{\footnotesize comm-H} & \fitfig{v1/comm-H}\\[7pt]
{\footnotesize comm-S} & \fitfig{v1/comm-S}\\[7pt]
{\footnotesize comm-CZ} & \fitfig{v1/comm-CZ}\\[7pt]
\end{longtable}
\end{document}
